\label{sec:algorithm}

\subsection{Economic Character Vector}

For each census tract $t$, let $C_{i}(t)$ denote the count of jobs in
NAICS sector $i$ (corresponding to WAC column \texttt{CNS0}$i$) and
$C_{0}(t)$ the total job count (column \texttt{C000}) after
block-to-tract aggregation.

We define three scalar signals:
\begin{align}
  \mathrm{CI}(t) &= \frac{C_{07}(t) + C_{09}(t) + C_{10}(t) + C_{11}(t)}{C_{0}(t)}
    \label{eq:ci} \\
  \mathrm{IF}(t) &= \frac{C_{01}(t) + C_{02}(t) + C_{05}(t) + C_{08}(t)}{C_{0}(t)}
    \label{eq:if} \\
  \mathrm{JPR}(t) &= \min\!\left(\frac{C_{0}(t)}{P(t)},\ 10.0\right)
    \label{eq:jpr}
\end{align}
where $P(t)$ is the total population of tract $t$ from the 2020 decennial
census, and the cap of 10.0 in equation~(\ref{eq:jpr}) prevents
pathologically high values in tracts with extremely high job concentration
and very low residential population (e.g., airport tarmacs, data centers).

\paragraph{Sector groupings.}
Equation~(\ref{eq:ci}) captures \emph{commercial intensity}: the fraction
of jobs in retail trade (CNS07), information (CNS09), finance and insurance
(CNS10), and real estate (CNS11).
These are the sectors characteristic of commercial corridors, downtown
business districts, and retail strips.

Equation~(\ref{eq:if}) captures \emph{industrial fraction}: the fraction
of jobs in agriculture and forestry (CNS01), mining and oil extraction
(CNS02), manufacturing (CNS05), and transportation and warehousing (CNS08).
These are the sectors characteristic of industrial parks, logistics hubs,
and extractive industry zones.

Equation~(\ref{eq:jpr}) captures \emph{jobs per resident}: the ratio of
total workplace jobs to resident population, capped at 10.
Values near zero indicate bedroom suburbs or rural residential areas.
Values above 1.0 indicate employment centers (more workers commute in than
residents).
Values above 5.0 indicate major employment concentrations such as university
research parks, hospital campuses, or industrial zones.

The economic character vector is
\begin{equation}
  \mathbf{e}(t) = \bigl(\mathrm{CI}(t),\ \mathrm{IF}(t),\ \mathrm{JPR}(t)\bigr)
  \in [0,1] \times [0,1] \times [0, 10].
  \label{eq:econ-char}
\end{equation}

For tracts with $C_0(t) = 0$ (no workplace jobs---purely residential tracts
with no employer establishments), we assign
$\mathbf{e}(t) = (0, 0, 0)$.
The cosine similarity is undefined for zero vectors and is handled
by the special cases in Section~\ref{subsec:cosine}.

\subsection{Cosine Similarity}
\label{subsec:cosine}

Given two adjacent tracts $u$ and $v$, their economic similarity is:
\begin{equation}
  \mathrm{sim}(\mathbf{e}_u, \mathbf{e}_v)
  \;=\;
  \begin{cases}
    1.0 & \text{if } \|\mathbf{e}_u\| = 0 \text{ and } \|\mathbf{e}_v\| = 0 \\
    0.5 & \text{if exactly one of } \|\mathbf{e}_u\|, \|\mathbf{e}_v\| = 0 \\
    \displaystyle\frac{\mathbf{e}_u \cdot \mathbf{e}_v}{\|\mathbf{e}_u\|\,\|\mathbf{e}_v\|}
        & \text{otherwise}
  \end{cases}
  \label{eq:cosine}
\end{equation}

The three cases handle:
\begin{itemize}
  \item \textbf{Both residential (zero-zero).}
    Two tracts with no workplace jobs are both purely residential.
    They should cluster together---similarity 1.0.
    This is the most common case in rural and suburban areas.
  \item \textbf{One residential, one employment (zero-nonzero).}
    A purely residential tract is adjacent to an employment center.
    The similarity is 0.5 (neutral): we neither strongly encourage nor
    strongly discourage this cut.
    The geographic boundary weight dominates in this case.
  \item \textbf{Both employment (nonzero-nonzero).}
    Standard cosine similarity applies.
    Two commercial corridors get similarity near 1.0;
    a commercial corridor next to an industrial park gets similarity
    near 0.1--0.3 (invite cut).
\end{itemize}

Mathematical properties:
\begin{itemize}
  \item \textbf{Symmetry.} $\mathrm{sim}(\mathbf{e}_u, \mathbf{e}_v)
    = \mathrm{sim}(\mathbf{e}_v, \mathbf{e}_u)$ for all pairs,
    satisfying METIS's requirement for symmetric edge weights.
  \item \textbf{Range.} $\mathrm{sim} \in [0, 1]$ always.
    Because all components of $\mathbf{e}$ are non-negative,
    the dot product is non-negative and cosine similarity cannot be negative.
  \item \textbf{Self-similarity.}
    $\mathrm{sim}(\mathbf{e}_t, \mathbf{e}_t) = 1.0$ for all tracts
    (a tract is maximally similar to itself).
\end{itemize}

\subsection{Weight Blending Formula}

Let $w_{\mathrm{geo}}(u, v)$ denote the existing geographic boundary-length
edge weight between adjacent tracts $u$ and $v$.
The blended economic character weight is:
\begin{equation}
  w_{\mathrm{ec}}(u, v)
  \;=\;
  w_{\mathrm{geo}}(u, v)
  \;\times\;
  \bigl(\alpha + (1 - \alpha)\,\mathrm{sim}(\mathbf{e}_u, \mathbf{e}_v)\bigr)
  \label{eq:blend}
\end{equation}
where $\alpha \in [0, 1]$ is the blend factor (default $\alpha = 0.5$).

At $\alpha = 0.5$:
\begin{itemize}
  \item Similar tracts ($\mathrm{sim} = 1.0$):
    $w_{\mathrm{ec}} = w_{\mathrm{geo}} \times 1.0$ --- weight unchanged.
  \item Dissimilar tracts ($\mathrm{sim} = 0.0$):
    $w_{\mathrm{ec}} = w_{\mathrm{geo}} \times 0.5$ --- weight halved.
\end{itemize}
The formula preserves the relative ordering of edge weights for pairs with
the same similarity: if $w_{\mathrm{geo}}(u,v) > w_{\mathrm{geo}}(x,y)$
and $\mathrm{sim}(\mathbf{e}_u, \mathbf{e}_v) = \mathrm{sim}(\mathbf{e}_x, \mathbf{e}_y)$,
then $w_{\mathrm{ec}}(u,v) > w_{\mathrm{ec}}(x,y)$.
This means geographic boundary length is never overridden by economic
similarity---it is modulated.
The $\alpha$ parameter controls the strength of the economic signal;
$\alpha = 1.0$ recovers pure geographic weights, $\alpha = 0.0$ makes
the geographic weight purely a scale factor multiplied by the similarity.
The tunable CLI parameter is \texttt{--econ-alpha}.

\subsection{Implementation}

The implementation lives in two new source files in the \texttt{bisect-cli} crate:

\paragraph{\texttt{lodes.rs}.}
Implements \texttt{load\_lodes\_wac\_tract(state, year)} which reads
\texttt{data/\{year\}/lodes/\{state\}\_wac\_tract.csv} and returns a
\texttt{HashMap<Geoid, EconChar>}.
If the file is absent, an empty map is returned (not an error), causing all
tracts to receive zero character vectors; the economic character weighter
then degrades gracefully to pure geographic weights (since all pairs get
$\mathrm{sim} = 1.0$, the zero-zero case).

The \texttt{EconChar} struct holds the three computed signals:
\begin{verbatim}
pub struct EconChar {
    pub commercial_intensity: f64,
    pub industrial_fraction:  f64,
    pub jobs_per_resident:    f64,
}
\end{verbatim}

\paragraph{\texttt{edge\_weights.rs} extension.}
The \texttt{EconomicCharacterWeighter} struct implements the \texttt{EdgeWeighter}
trait, applying equation~(\ref{eq:blend}) to each edge in the weight map:
\begin{verbatim}
let sim = cosine_similarity(
    chars.get(&u).unwrap_or(&ZERO_CHAR),
    chars.get(&v).unwrap_or(&ZERO_CHAR),
);
let w_new = alpha * w + (1.0 - alpha) * sim * w;
\end{verbatim}
This step is $O(|E|)$, where $|E|$ is the number of adjacency-graph edges
(approximately $3n$ for typical planar tract graphs).
On a 2000-tract state, the computation adds under 50\,ms of wall time.

\paragraph{CLI wiring.}
Selecting \texttt{--weights-override economic-character} activates
\texttt{WeightMode::EconomicCharacter}, which sets \texttt{WeightSpec\{geographic: true, economic\_character: true, econ\_alpha: 0.5\}}.
The LODES data requirement is satisfied by running
\texttt{bisect fetch --type lodes --year 2020 --state \{state\}} before the plan build.

\paragraph{L0 test invariants.}
The implementation is verified by the following unit test invariants checked
on every build:
\begin{itemize}
  \item \texttt{cosine\_sim\_both\_zero\_returns\_one}
  \item \texttt{cosine\_sim\_one\_zero\_returns\_half}
  \item \texttt{cosine\_sim\_identical\_returns\_one}
  \item \texttt{cosine\_sim\_orthogonal\_returns\_zero}
  \item \texttt{econ\_weighter\_leaves\_similar\_tracts\_unchanged}
    ($\alpha=0.5$, $\mathrm{sim}=1.0 \Rightarrow$ weight unchanged)
  \item \texttt{econ\_weighter\_halves\_dissimilar\_edges}
    ($\alpha=0.5$, $\mathrm{sim}=0.0 \Rightarrow$ weight halved)
  \item \texttt{load\_lodes\_missing\_file\_returns\_empty}
  \item \texttt{align\_lodes\_to\_adjacency\_zero\_for\_missing}
\end{itemize}
